\documentclass[10pt,twocolumn]{article}
\usepackage[T1]{fontenc}
\usepackage[utf8]{inputenc}
\usepackage{lmodern}
\usepackage[margin=1.8cm,top=2.4cm,bottom=2cm,columnsep=0.6cm]{geometry}
\usepackage{fancyhdr}
\usepackage{titlesec}
\usepackage{booktabs}
\usepackage{amsmath,amssymb,amsfonts}
\usepackage{microtype}
\usepackage{xcolor}
\usepackage{mdframed}
\usepackage{parskip}
\usepackage{enumitem}
\usepackage{hyperref}

\definecolor{rulered}{RGB}{180,30,30}
\definecolor{boxgray}{RGB}{245,245,245}
\definecolor{keywordgray}{RGB}{100,100,100}

\pagestyle{fancy}
\fancyhf{}
\fancyhead[L]{\textsc{\small Claude Mag}}
\fancyhead[C]{\small\textit{Machine Learning Research Digest}}
\fancyhead[R]{\small 2026-11-16}
\fancyfoot[C]{\small\thepage}
\renewcommand{\headrulewidth}{0.8pt}

\titleformat{\section}{\normalsize\bfseries\color{rulered}}{}{0pt}{}%
  [\vspace{-4pt}\color{rulered}\rule{\columnwidth}{0.4pt}]
\titlespacing{\section}{0pt}{8pt}{4pt}

\hypersetup{colorlinks=true,urlcolor=rulered,linkcolor=black}

\begin{document}
\setlength{\parindent}{0pt}
\setlength{\parskip}{4pt}

{\small\color{keywordgray}\textbf{Keywords:}
  Drug Discovery \textbullet{} Molecular Generation \textbullet{} Reinforcement Learning \textbullet{} Protein Binding \textbullet{} NeurIPS 2026}\par\vspace{4pt}

{\LARGE\bfseries\raggedright Binding Mode Matters: Hotspot-Aware Drug Discovery via Explorative Preferences}\par\vspace{4pt}
{\large\itshape\raggedright Decomposing the Protein Pocket into Hotspots Unlocks Structurally Diverse, High-Affinity Compounds}\par\vspace{6pt}
{\small\textbf{By} Dingshuo Chen, Qian Zhao, Xiangxiang Zeng, et al.}\par\vspace{2pt}
{\small\color{keywordgray}arXiv:2611.14956 \textbullet{} cs.LG, q-bio.BM \textbullet{} NeurIPS 2026}
\par\vspace{2pt}\noindent{\color{rulered}\rule{\columnwidth}{0.6pt}}\vspace{6pt}

A single docking score hides the multifaceted nature of protein-ligand binding and drives AI drug discovery into mode collapse---generating many similar molecules that all exploit the same sub-region of the binding pocket.
Chen~et~al.\ introduce BindMol, a reinforcement learning framework that decomposes the pocket into binding hotspots and assigns a separate reward to each, steering molecular generation toward structurally diverse compounds that engage the target through multiple distinct binding modes.
Accepted at NeurIPS 2026.

\begin{mdframed}[backgroundcolor=boxgray,linecolor=rulered,linewidth=1pt,
    innerleftmargin=6pt,innerrightmargin=6pt,
    innertopmargin=6pt,innerbottommargin=6pt]
\textbf{\small AT A GLANCE}\par\vspace{4pt}
\begin{description}[leftmargin=0pt,labelsep=4pt,itemsep=2pt,parsep=0pt,
    font=\small\bfseries]
  \item[Problem:] Single-score optimisation collapses molecular generation to one chemotype and misses structurally distinct binding modes.
  \item[Why it matters:] Diverse binding modes yield diverse chemical scaffolds, providing backup leads and improving drug campaign success rates.
  \item[Method:] Cluster pocket residues into $K$ hotspots; assign per-hotspot RL rewards with an exploration bonus; fine-tune a 3D autoregressive generator.
  \item[Result:] Average Vina score $-9.4$ kcal/mol vs.\ $-8.7$ for best single-score baseline; scaffold diversity 0.72 vs.\ 0.41.
  \item[Evidence:] CrossDocked2020 benchmark plus three additional PDB targets; ablations over $K$, exploration bonus, and reward weighting.
\end{description}
\end{mdframed}

\section{The Big Picture}

Structure-based drug design (SBDD) uses the 3D structure of a protein target to guide the generation of molecules that bind tightly.
The standard computational workflow optimises a docking score---a scalar estimate of binding free energy---as the sole objective.
Generative models trained or fine-tuned with this signal learn to find molecules that score well, but \emph{scoring well} and \emph{binding diversely} are not the same.

A binding pocket is not a single site: it is a landscape of sub-regions (hotspots) with distinct chemical characters---hydrophobic cores, hydrogen-bond donors, salt bridge anchors.
Different chemotypes can achieve high affinity by engaging different subsets of these hotspots.
Optimising a single composite score obliterates this information.

BindMol reinstates the multi-modal nature of binding by making hotspot diversity a first-class training objective.

\section{Why Existing Approaches Fall Short}

\textbf{Mode collapse in molecular generation.}
When a generative model is fine-tuned to maximise a single docking score, it quickly converges to exploiting one dominant hotspot---typically the most hydrophobic sub-pocket---and generates many structurally similar molecules that share the same scaffold.
This is mode collapse at the level of chemical space.

\textbf{Loss of chemical diversity.}
In a real drug discovery campaign, medicinal chemists need a diverse library to hedge against toxicity, selectivity, and synthetic accessibility issues.
A library of similar compounds offers few backup options.

\textbf{Failure to explore.}
Prior methods that add diversity objectives (e.g., Tanimoto similarity penalties) operate at the level of whole molecules, not binding interactions.
They promote molecular diversity without connecting it to functional differences in how molecules engage the target.

\section{Core Mechanism}

BindMol decomposes the standard SBDD problem into three components:

\begin{enumerate}[leftmargin=1.5em,itemsep=2pt]
  \item \textbf{Hotspot clustering:} Given a protein pocket and reference ligand, cluster pocket residues into $K$ hotspots.
  \item \textbf{Per-hotspot rewards:} Assign a separate reward to each hotspot measuring how much a generated molecule interacts with that hotspot.
  \item \textbf{Explorative preference:} Add a diversity bonus that rewards molecules whose per-hotspot score profiles differ from previously generated molecules, encouraging exploration of the binding mode space.
\end{enumerate}

The molecular generator is an autoregressive 3D graph model that places atoms one at a time within the pocket, guided by the multi-objective RL reward.

\section{Technical Formulation}

\subsection*{Hotspot Clustering}

Let the pocket be defined by the set of residues $\mathcal{P} = \{p_1, \ldots, p_M\}$ within 10\,\AA{} of the reference ligand.
Hotspots are defined by hierarchical agglomerative clustering on the residue feature matrix:

\begin{equation}
F(p_i) = \bigl[\mathbf{c}(p_i);\; \psi(p_i);\; \mathrm{SASA}(p_i)\bigr]
\end{equation}

where $\mathbf{c}(p_i) \in \mathbb{R}^3$ is the $C_\alpha$ coordinate, $\psi(p_i)$ is a physicochemical feature vector (hydrophobicity, hydrogen-bond capacity, charge), and $\mathrm{SASA}(p_i)$ is the solvent-accessible surface area.
Clustering with linkage distance threshold $d_c$ produces $K$ hotspot clusters $H_1, \ldots, H_K$.

\subsection*{Per-Hotspot Reward}

For a generated molecule $m$ docked into the pocket with binding pose $\mathcal{P}(m)$:

\begin{equation}
r_k(m) = \sum_{i \in H_k} w_i \cdot I(m, p_i)
\end{equation}

where $I(m,p_i)$ is the residue-specific interaction energy from the docking software and $w_i$ is a weight proportional to the residue's contribution to the reference ligand's binding energy.

\subsection*{Aggregate Multi-Objective Reward}

The total reward is:

\begin{equation}
\mathcal{R}(m) = \sum_{k=1}^K \bar{w}_k \cdot r_k(m) + \lambda \cdot D(m;\, \mathcal{M}_{t})
\label{eq:reward}
\end{equation}

where $\bar{w}_k = 1/K$ (uniform by default), $\mathcal{M}_{t}$ is the set of molecules generated so far in the current episode, and $D$ is the diversity bonus:

\begin{equation}
D(m;\, \mathcal{M}_t) = \frac{1}{|\mathcal{M}_t|} \sum_{m' \in \mathcal{M}_t} \bigl\|\mathbf{r}(m) - \mathbf{r}(m')\bigr\|_2
\end{equation}

This penalises molecules whose hotspot score \emph{profile} $\mathbf{r}(m) = (r_1(m),\ldots,r_K(m))$ is similar to previously generated molecules, pushing the generator to explore different binding mode signatures.

\subsection*{Policy Gradient Training}

The generative model $p_\theta(m \mid \text{pocket})$ is fine-tuned with REINFORCE:

\begin{equation}
\nabla_\theta \mathcal{J}(\theta) = \mathbb{E}_{m \sim p_\theta}\bigl[\mathcal{R}(m) \cdot \nabla_\theta \log p_\theta(m \mid \text{pocket})\bigr]
\end{equation}

with a KL regularisation term towards the pre-trained generator to prevent reward hacking:

\begin{equation}
\mathcal{J}_{\mathrm{reg}}(\theta) = \mathcal{J}(\theta) - \beta_{\mathrm{KL}} \cdot D_{\mathrm{KL}}\bigl(p_\theta(\cdot \mid \text{pocket}) \,\|\, p_{\theta_0}(\cdot \mid \text{pocket})\bigr)
\end{equation}

\section{Learning or Inference Procedure}

\textbf{Pre-training.}
The 3D autoregressive generator is pre-trained on CrossDocked2020 (22M docked pose pairs) using maximum likelihood estimation to generate valid, drug-like molecules in protein pockets.

\textbf{RL fine-tuning.}
For each target, run 10K rollouts of molecule generation, compute per-hotspot rewards via a fast docking oracle (Vina), and update $\theta$ via policy gradient with KL regularisation.
Hotspot clustering is computed once per target and cached.
Training takes approximately 4 GPU-hours per target.

\textbf{Generation.}
At inference, sample $M=100$ molecules per target and filter by drug-likeness (QED $>$ 0.5) and synthetic accessibility (SA score $<$ 3).

\section{What the Guarantee Says}

Under the assumption that the docking score is a linear function of per-hotspot interaction scores plus noise, the optimal policy under Eq.~(\ref{eq:reward}) at $\lambda=0$ is equivalent to maximising the sum of per-hotspot scores, which equals maximising the original docking score.
Adding $\lambda > 0$ creates a Pareto expansion: the optimal policy explores the full Pareto front of per-hotspot scores, recovering the single-score optimum as a special case ($\lambda \to 0$) and maximising diversity as another ($\lambda \to \infty$).
The choice $\lambda = 0.3$ is validated by cross-validation on held-out targets.

\section{Experimental Findings}

\begin{center}
\small
\begin{tabular}{lrrrr}
\toprule
\textbf{Method} & \textbf{Vina} & \textbf{Div.} & \textbf{QED} & \textbf{SA$<$3} \\
\midrule
AR-baseline & $-7.9$ & 0.31 & 0.68 & 76\% \\
DiffSBDD & $-8.3$ & 0.37 & 0.70 & 78\% \\
TargetDiff & $-8.7$ & 0.41 & 0.68 & 79\% \\
BindMol (ours) & $\mathbf{-9.4}$ & $\mathbf{0.72}$ & $0.72$ & $83\%$ \\
\bottomrule
\end{tabular}
\end{center}

Vina scores are in kcal/mol (more negative is better); Div.\ is Bemis-Murcko scaffold diversity (Tanimoto-based); QED is quantitative estimate of drug-likeness; SA$<$3 is fraction of molecules with synthetic accessibility score below 3.

BindMol achieves the best Vina score \emph{and} the highest scaffold diversity simultaneously---a combination no single-score baseline achieves.

\section{Ablations and Interpretation}

\begin{itemize}[leftmargin=1em,itemsep=2pt]
  \item \textbf{$K=2$ vs.\ $K=4$:} $K=4$ improves diversity on large flexible pockets (CDK2, HSP90) but adds no benefit on compact rigid pockets (thrombin). Optimal $K$ correlates with pocket volume.
  \item \textbf{Diversity bonus $\lambda$:} $\lambda=0$ (no exploration) recovers near-TargetDiff performance. $\lambda=1.0$ over-diversifies and hurts Vina score. $\lambda=0.3$ is Pareto-optimal.
  \item \textbf{Profile diversity vs.\ molecular diversity:} Replacing the hotspot-profile diversity bonus with Tanimoto molecular diversity reduces Vina score by 0.4 kcal/mol---confirming that diversity at the binding-mode level is more informative than chemical-fingerprint diversity.
\end{itemize}

\section*{Reference}
D.\ Chen, Q.\ Zhao, X.\ Zeng, et al.
\textbf{Binding Mode Matters: Hotspot-Aware Drug Discovery via Explorative Preferences.}
arXiv:2611.14956. cs.LG, q-bio.BM. NeurIPS 2026.
\url{https://arxiv.org/abs/2611.14956}

\end{document}
